% SPDX-License-Identifier: CC-BY-SA-4.0
% Author: Matthieu Perrin

\begingroup

\begin{exercice}[Mini-défi -- Intersection de langages rationnels]\label{exo:lexing/kleene/challenge}

  On travaille sur l’alphabet $\Sigma=\{a,b\}$. On s’intéresse au langage
  $$
  L \eqdef \{ u\in\{a,b\}^\star \mid |u|_a \equiv 0 \pmod 2 \land |u|_b \equiv 0 \pmod 3 \}.
  $$

  \begin{question}
  \item Construire un automate fini déterministe $A$ qui reconnaît les mots contenant un nombre pair de $a$.
  \item Construire un automate fini déterministe $B$ qui reconnaît les mots contenant un nombre multiple de $3$ de $b$.
  \item Construire le produit cartésien $A \times B$ et préciser l'ensemble des états acceptants.
  \item Poser le système d'équations rationnelles reliant les langages droits des états de $A \times B$.
  \item En déduire une expression rationnelle pour $L$.
  \end{question}

\end{exercice}

\begin{correction}

  \begin{question}

  \item Deux états selon la parité du nombre de $a$ lus ; $b$ ne change pas la parité.
    \begin{tikzpicture}[automaton, baseline=(E.base)]
      \state[initial, accepting] (E) at (0,1) {E};
      \state                     (O) at (0,0) {O};
      \path (E) edge[loop right] node       {$b$} (E);
      \path (O) edge[loop right] node       {$b$} (O);
      \path (E) edge[bend right] node[swap] {$a$} (O);
      \path (O) edge[bend right] node[swap] {$a$} (E);
    \end{tikzpicture}

  \item
    \begin{tikzpicture}[automaton, baseline=(R0.base)]
      \state[initial, accepting] (R0) at (0,0) {$0$};
      \state                     (R1) at (1,0) {$1$};
      \state                     (R2) at (2,0) {$2$};
      \path (R0) edge[loop above] node {$a$} (R0);
      \path (R1) edge[loop above] node {$a$} (R1);
      \path (R2) edge[loop above] node {$a$} (R2);
      \path (R0) edge             node {$b$} (R1);
      \path (R1) edge             node {$b$} (R2);
      \path (R2) edge[bend left]  node {$b$} (R0);
    \end{tikzpicture}

  \item
    \begin{tikzpicture}[automaton, baseline=(E0.base), x=20mm, y=20mm]
      \state[initial, accepting] (E0) at (0,1) {$E,0$};
      \state                     (O0) at (0,0) {$O,0$};
      \state                     (E1) at (1,1) {$E,1$};
      \state                     (O1) at (1,0) {$O,1$};
      \state                     (E2) at (2,1) {$E,2$};
      \state                     (O2) at (2,0) {$O,2$};

      \path (E0) edge[bend left]  node       {$a$} (O0);
      \path (O0) edge[bend left]  node       {$a$} (E0);
      \path (E1) edge[bend left]  node       {$a$} (O1);
      \path (O1) edge[bend left]  node       {$a$} (E1);
      \path (E2) edge[bend left]  node       {$a$} (O2);
      \path (O2) edge[bend left]  node       {$a$} (E2);

      \path (E0) edge             node[swap] {$b$} (E1);
      \path (E1) edge             node[swap] {$b$} (E2);
      \path (E2) edge[bend right] node[swap] {$b$} (E0);
      \path (O0) edge             node       {$b$} (O1);
      \path (O1) edge             node       {$b$} (O2);
      \path (O2) edge[bend left]  node       {$b$} (O0);
    \end{tikzpicture}

  \item En notant $L_{E0}$, $L_{O0}$, $L_{E1}$, $L_{O1}$, $L_{E2}$, $L_{O2}$
    les langages droits des états $(E,0)$, $(O,0)$, $(E,1)$, $(O,1)$, $(E,2)$, $(O,2)$ :
    $$
    \left\{
    \begin{array}{rcl}
      L_{E0} &=& aL_{O0} \mid bL_{E1} \mid \varepsilon\\
      L_{O0} &=& aL_{E0} \mid bL_{O1}\\
      L_{E1} &=& aL_{O1} \mid bL_{E2}\\
      L_{O1} &=& aL_{E1} \mid bL_{O2}\\
      L_{E2} &=& aL_{O2} \mid bL_{E0}\\
      L_{O2} &=& aL_{E2} \mid bL_{O0}
    \end{array}
    \right.
    $$

  \item Une expression rationnelle pour $L$ (langage droit de l’état initial $(E,0)$) :
    $$
    L = \Big( \big((tU)^\star\, tV\, Q^\star\, R\big)^\star\, (tU)^\star\, t \Big)
    $$
    avec
    \[
    t \eqdef \big(a(bbb)^\star a\big)^\star
    \]
    \[
    Q \eqdef t\,(b \mid ab(bbb)^\star a)\,t\,a(bbb)^\star b^2 a
    \]
    \[
    R \eqdef t\Big((b \mid ab(bbb)^\star a)\,t\,(b \mid a(bbb)^\star b a)\ \mid\ ab^2(bbb)^\star a\Big)
    \]
    \[
    U \eqdef ab^2(bbb)^\star a\;t\,(b \mid a(bbb)^\star b a)
    \]
    \[
    V \eqdef (b \mid aba \mid ab^4(bbb)^\star a)\ \mid\ ab^2(bbb)^\star a\;t\;a(bbb)^\star b^2 a
    \]
    
  \end{question}

\end{correction}

\endgroup
\endinput
